\begin{turkishabstract}\label{chp:turkishabstract}
\vspace{-1.5em}

Büyük ölçekli GPU kümelerindeki temel zorluk, zamanlayıcıların paylaşılan donanıma atanan heterojen, ağır kuyruklu işleri öngörememesi veya kontrol edememesidir. İş yüklerinin uzun kuyruklu doğası nedeniyle, devasa bir ``fil'' işinin kendinden daha küçük birçok ``fare'' görevini engellemesiyle GPU kümelerinde aşırı kuyruk başı tıkanması (head-of-line blocking) yaşanır. Geleneksel Önce Tahmin Et Sonra Optimize Et (Predict-then-Optimize - PtO) çerçevesi, zamanlamadan önce nokta tahmin hatalarını (örneğin, ortalama mutlak hata) en aza indirerek bunu çözmeye çalışır. Bu tezde, makine öğrenmesine yeni bir yön kazandırmak amacıyla Karar Odaklı Öğrenme (Decision-Focused Learning - DFL) yaklaşımını öneriyoruz. Nihayetinde tahminleyicinin amacının yalnızca kendi başına doğru tahminler elde etmek değil, genel sistem düzeyindeki performans ölçütlerini, spesifik olarak da İş Tamamlanma Süresini (JCT) optimize etmek olduğunu gösteriyoruz. Alibaba'nın Yapay Zeka Platformundan (PAI) alınan verileri kullanarak, çalışma süresi tahminlerini hem küçük (32-GPU) hem de devasa (256-GPU) küme ortamlarında olay tahrikli (event-driven) simülasyonlara bağlayan uçtan uca bir test ortamı geliştirdik. Altı modelimizin tamamı, gönderim anında hesaplanan, sweep-line tabanlı yeni bir küme-kullanımı özniteliğini girdi olarak kullanmaktadır. Analiz sonuçları, bir modelin nokta-tahmin doğruluğunun (ortalama mutlak hata) hangi politikanın en büyük zamanlama kazancını sağlayacağını güvenilir biçimde göstermediğini ortaya koymaktadır. Kategorik LSTM politikası, simülasyonlarımızda hem 32-GPU hem de 256-GPU ölçeğinde -- daha yüksek nokta-tahmin hatasına rağmen tüm ağaç tabanlı modellerin önünde -- FIFO'ya kıyasla en yüksek JCT azalmasını sağlamıştır; değerlendirilen tüm politikaların iş-tamamlanma-süresi sıralaması da iki küme boyutu arasında neredeyse birebir aynıdır (Spearman sıra korelasyonu 0,98'in üzerinde).\footnote{Kesin JCT ve sıra-korelasyonu rakamları, hiperparametre ayarlama kayıp fonksiyonundaki kod düzeyi düzeltmenin ve öğrenilmemiş referans taban çizgilerinin (Kullanıcı-Medyanı, Profil-Medyanı) eklenmesinin ardından yenilenmektedir; bkz. Sınırlamalar bölümü.} Derin öğrenme modelleri daha yüksek nokta hataları üretse de, ``fareleri'' ``fillerden'' koruyan üstün birer ``Sıralama Öğrenmesi'' (Learning to Rank) motoru gibi hareket etmektedirler. Bu sonuçlar, doğru göreceli sıralama kararlarını vermenin doğru mutlak tahmini elde etmekten daha önemli olduğuna dair somut kanıtlar sunmaktadır -- ve mevcut bulgulara göre bu, yalnızca test edilen en büyük ölçekte değil, küme boyutundan bağımsız olarak geçerlidir.

\end{turkishabstract}